%=============================================================================!
% 3.5  SLIDING ON A HILL-SHAPED TRACK                                         !
%=============================================================================!
\section{Sliding on a hill-shaped track}
\label{sec:en-track}

A ball thrown upwards moves along a straight line. Many motions follow a
curved track instead, as a sledge follows the slopes of a hillside, and
energy is most useful there, because the forces along a curved track
change from point to point. This section treats an object sliding along a
track shaped into hills and valleys, and finds its speed at every point
from its height alone.

\subsection*{A bead on a wire}

The cleanest version of such a track is a smooth, stiff wire bent into
hills and valleys in a vertical plane, with a bead threaded on it. We
measure $x$ horizontally and $y$ vertically upwards, and describe the
shape of the wire by its height $h(x)$ above the ground; our example,
drawn in \cref{fig:en-track}(a) in the next section, is
\begin{equation}
   h(x) = 0.15\,x^4 - 0.6\,x^2 + 0.15\,x + 1.0 ,
   \label{eq:en-track}
\end{equation}
with $x$ and $h$ in metres. It has a deep valley on the left, a hill in
the middle and a shallower valley on the right. Three features of this
system are chosen so that it obeys the simple laws below, and each is an
idealisation that a real track only approaches.
\begin{itemize}
\item The bead is threaded on the wire, so it cannot leave the track. A
   block sliding on top of a track could lose contact over a crest, where
   the track curves downwards more sharply than the fast-moving block can
   follow, and then fly through the air instead.
\item The bead slides. A ball rolling along a track also spins, and part
   of its kinetic energy is then in the spinning, so at each height a
   rolling ball moves more slowly than a sliding body; Chapter~5 treats
   spinning.
\item The wire is smooth. No real wire is perfectly smooth: friction
   slowly removes energy (\cref{sec:en-friction}), and a real bead never
   quite returns to the height from which it was released.
\end{itemize}

\subsection*{The push of the wire does no work}

Two forces act on the bead: its weight, $m\vb{g}$ downwards, and the
push of the wire. A smooth surface exerts no friction, so the push of the
wire has no component along the wire: it is a normal force,
perpendicular to the wire at the point where the bead is
(\cref{fig:en-bead}a). Its size and direction change as the bead moves,
and we do not need to know them. The velocity of the bead points along
its path (\cref{sec:mo-plane}), which is along the wire, so the normal
force $\vb{F}_{\mathrm{n}}$ is perpendicular to the velocity at every
instant, $\vb{F}_{\mathrm{n}}\cdot\vb{v} = 0$. Its power is therefore
zero throughout, so the work it does over any part of the motion, the
integral of its power as in \cref{eq:en-work-energy}, is zero.

\begin{figure}[htb]
\centering
\begin{tikzpicture}[line cap=round, line join=round, scale=1.3,
   force/.style={-{Stealth[length=5pt]}, line width=1.0pt}]
   % (a) the forces on the bead
   \begin{scope}
      \node at (0.1,2.25) {\small (a)};
      \draw[line width=0.9pt, black!70, domain=0:3.6, smooth, samples=60]
         plot (\x, {1.1 + 0.55*cos(deg(1.3*\x)) + 0.05*\x});
      \coordinate (P) at (1.6,0.912);
      \draw[force] (P) -- ++(0.498*0.95, 0.867*0.95)
         node[above] {\small $\vb{F}_{\mathrm{n}}$};
      \draw[force] (P) -- ++(0,-0.95) node[below] {\small $m\vb{g}$};
      \draw[force, accent, dashed] (P) -- ++(0.867*0.9, -0.498*0.9)
         node[right] {\small $\vb{v}$};
      \fill[black!15, draw=black, line width=0.5pt] (P) circle (0.09);
   \end{scope}
   % (b) the wire cut into short chords
   \begin{scope}[xshift=4.3cm]
      \node at (0.1,2.25) {\small (b)};
      \draw[line width=0.5pt, black!35, domain=0:5.0, smooth, samples=60]
         plot (\x, {1.1 + 0.55*cos(deg(1.3*\x)) + 0.05*\x});
      \draw[line width=0.8pt, accent] (0.2,1.642) -- (1.1,1.232)
         -- (2.0,0.729) -- (2.9,0.800) -- (3.8,1.414) -- (4.7,1.877);
      \foreach \x/\y in {0.2/1.642, 1.1/1.232, 2.0/0.729, 2.9/0.800,
         3.8/1.414, 4.7/1.877}
         \fill[accent] (\x,\y) circle (1.3pt);
      \draw[black!60, dashed, line width=0.4pt] (2.9,0.800) -- (3.8,0.800);
      \draw[{Stealth[length=3pt]}-{Stealth[length=3pt]}, line width=0.4pt]
         (2.9,0.62) -- (3.8,0.62) node[midway, below] {\small $\Delta x_k$};
      \draw[{Stealth[length=3pt]}-{Stealth[length=3pt]}, line width=0.4pt]
         (3.95,0.800) -- (3.95,1.414) node[midway, right] {\small $\Delta y_k$};
      \node[below] at (0.2,1.58) {\small $A$};
      \node[above] at (4.7,1.95) {\small $B$};
   \end{scope}
\end{tikzpicture}
\caption{A bead on a smooth wire. (a) The two forces on the bead: its
weight $m\vb{g}$ and the push $\vb{F}_{\mathrm{n}}$ of the wire, which is
perpendicular to the wire and so to the velocity $\vb{v}$. (b) The wire
from $A$ to $B$ replaced by short straight chords; on the chord that
rises by $\Delta y_k$ the weight does the work $-mg\,\Delta y_k$.}
\label{fig:en-bead}
\end{figure}

\subsection*{The work of the weight}

The weight is constant, but the path is curved. To extend the straight-path
argument of \cref{eq:en-work-constant}, we replace the wire between two
points $A$ and $B$ by a chain of short straight chords
(\cref{fig:en-bead}b). On the $k$th chord the displacement is $(\Delta
x_k, \Delta y_k)$, and the weight, $(0, -mg)$, does the work
\begin{equation*}
   (0, -mg)\cdot(\Delta x_k, \Delta y_k) = -mg\,\Delta y_k .
\end{equation*}
Adding the chords, the rises $\Delta y_k$ add up to the total rise from
$A$ to $B$, since each chord starts where the one before it ends, so
\begin{equation*}
   W_{\mathrm{weight}} = -mg\sum_k\Delta y_k = -mg\,(y_B - y_A) .
\end{equation*}
The result is the same for any number of chords, so it holds also in the
limit in which the chords follow the wire exactly. The weight does work
that depends only on the change in height, whatever the shape of the path
between.

\subsection*{The speed at every height}

By the work-energy theorem \cref{eq:en-work-energy}, the change in the
bead's kinetic energy is the work of the weight plus the work of the
wire, $-mg\,(y_B - y_A) + 0$. On the wire $y = h(x)$, so
\begin{align*}
   \tfrac12 m v_B^2 - \tfrac12 m v_A^2
   &= -mg\,\bigl(h(x_B) - h(x_A)\bigr)\\
   &= -mg\,h(x_B) + mg\,h(x_A)
      && \text{multiplying out the bracket,}\\
   \tfrac12 m v_B^2 + mg\,h(x_B) &= \tfrac12 m v_A^2 + mg\,h(x_A)
      && \text{adding $mg\,h(x_B) + \tfrac12 m v_A^2$ to both sides.}
\end{align*}
The total energy
\begin{equation}
   E = \tfrac12 m v^2 + mg\,h(x)
   \label{eq:en-track-energy}
\end{equation}
is conserved, with the potential energy $U = mgh(x)$ of the weight. It is
convenient to write $E = mgh_E$, so that $h_E$ is the height at which the bead
would be at rest; a bead released from rest at height $h_E$ has exactly
this energy. Then $\frac12 mv^2 = mg\,(h_E - h)$, and dividing by $m/2$ and
taking the square root,
\begin{equation*}
   v = \sqrt{2g\,\bigl(h_E - h(x)\bigr)} .
\end{equation*}
The mass has cancelled: every bead released from the same height moves
at the same speed at the same place. Released from rest at $h_E =
\SI{1.3}{\metre}$, the bead passes the bottom of the deep valley, where
$h = \SI{0.1834}{\metre}$, at $\sqrt{2\times9.80665\times(1.3 - 0.1834)} =
\SI{4.680}{\metre\per\second}$, crosses the top of the hill, at
\SI{1.0094}{\metre}, at \SI{2.387}{\metre\per\second}, and passes the
bottom of the shallow valley, at \SI{0.6072}{\metre}, at
\SI{3.686}{\metre\per\second}. None of these numbers required the shape
of the wire between the points, or the time taken to get there.

Energy gives the speed at each place directly. It does not say which way
the bead is moving: a bead passing a point to the left and one passing it
to the right have the same speed. Nor does it give the time of arrival at
once; that needs the time spent on each short stretch of wire, its length
divided by the speed there, added up along the way, or the equation of
motion along the wire, which Chapter~6 derives in a form suited to such
constrained motion.

\notebook{03}{release the bead from different heights and read its speed}

\begin{selfcheck}
A bead is released from rest on the left slope at a height of
\SI{0.8}{\metre}. Can it reach the shallow valley on the right? What is
the least speed it would need at the bottom of the deep valley to cross
the hill?
\end{selfcheck}
